\begin{abstract}
Retrieval-augmented language models increasingly depend on a pool of
heterogeneous candidate passages, and the correctness of a grounded
answer hinges on which passages enter the context window. Most
systems collapse this selection to a single relevance score, yet a
classical view of trust decomposes into several orthogonal
dimensions; on HotpotQA's distractor setting we can directly
measure three of them: relevance to the query, an intrinsic
authority proxy, and corroboration across independent candidate
passages. We ask whether combining six lightweight signals covering
these three dimensions yields better selection than any single
signal, and whether the textbook intuition for combining them is
empirically correct. On HotpotQA's distractor setting, with exactly two gold
paragraphs among ten candidates per question, we compare a naive
uniform average of three trust signals against a six-feature
logistic regression fit on only one hundred questions per seed. The
learned combiner attains Precision@2 of $0.345 \pm 0.032$, a 3.3
absolute-point improvement over the strongest single-signal baseline
SBERT at $0.312 \pm 0.035$, with the gap consistent on every seed;
the gold-versus-distractor score gap more than doubles, from
$0.170$ to $0.349$, indicating substantially better calibration. The
uniform average instead collapses to $0.080$, below BM25. Inspecting
the fitted weights diagnoses why: paragraph length and cross-
paragraph consistency are learned as anti-signals, because
distractors are length-matched to gold and topically close enough
that cross-source similarity peaks among distractors rather than
among the gold pair. In distractor-rich retrieval, the classical
trust heuristics are correct in prior but wrong in sign, and
supervised calibration over only one hundred questions per seed is
sufficient to recover them. Evidence comes from a single benchmark,
so the sign-flip should be read as a concrete existence proof rather
than a universal property of distractor-rich retrieval.
\end{abstract}
